\textbf{\textit{Soal. }} Diberikan persegi \(ABCD\) dan titik \(E\) pada ruas \(CD\). Misalkan \(F\) adalah titik pada ruas \(AE\) sehingga \(BF\) tegak lurus \(AE\). Dengan cara serupa, misalkan \(G\) adalah titik pada ruas \(EB\) sehingga \(AG\) tegak lurus \(EB\). Misalkan garis \(DF\) dan \(CG\) berpotongan di \(H\). Buktikan bahwa
\[
\angle DFC + \angle DGC = 90^\circ + \angle AHD + \angle BHC.
\]
\textit{Proposed by: Raymond Christopher Tanto}
\begin{proof}[\textbf{Solusi. }] (Minggu, 26 Januari 2025) 
    % \begin{center}
    \begin{asy}
        // Import the Asymptote geometry module
        import olympiad;
        import geometry;

        // Define the size of the square
        size(350);

        // Define the square ABCD
        point A = (0, 0);
        point B = (1, 0);
        point C = (1, 1);
        point D = (0, 1);

        // Define center
        point O = (0.5, 0.5);

        // Draw the square
        draw(A--B--C--D--cycle);

        // Define point E on segment CD
        point E = (0.7, 1);

        // Draw segment AE
        draw(A--E);

        // Construct point F on AE such that BF is perpendicular to AE
        line AE = line(A, E);
        line BF = perpendicular(B, AE);
        point F = intersectionpoint(AE, BF);

        // Draw segment BF
        draw(B--F);

        // Construct point G on EB such that AG is perpendicular to EB
        line EB = line(E, B);
        line AG = perpendicular(A, EB);
        point G = intersectionpoint(EB, AG);

        // Draw segment AG
        draw(A--G);

        // Draw lines DF and CG
        line DF = line(D, F);
        line CG = line(C, G);

        // Find the intersection point H of DF and CG
        point H = intersectionpoint(DF, CG);
        point T = intersectionpoint(AG, BF);

        // Draw lines DF and CG
        draw(D--F);
        draw(C--G);

        draw(E--B);
        draw(F--H--G);
        draw(A--H--B);

        draw(D--B, red);
        draw(A--C, red);
        draw(O--E, red);
        draw(O--F, red);

        draw(F--C, dotted+red);
        draw(G--D, dotted+red);

        // Label the points
        label("$A$", A, SW);
        label("$B$", B, SE);
        label("$C$", C, NE);
        label("$D$", D, NW);
        label("$E$", E, N);
        label("$F$", F, N);
        label("$G$", G, NE);
        label("$H$", H, S);
        label("$O$", O, S);
        label("$T$", T, S);

        // Mark right angles
        draw(rightanglemark(B, F, A, 1));
        draw(rightanglemark(A, G, E, 1));
        draw(rightanglemark(C, D, A, 1));
        draw(rightanglemark(B, C, D, 1));
        draw(rightanglemark(C, O, D, 1), red);

        // Add circles
        // Circle through A, B, G
        circle circleABG = circle(A, B, G);
        draw(circleABG, black);

        // Circle through D, F, E
        circle circleDFE = circle(D, F, E);
        draw(circleDFE, dashed + blue);

        // Circle through E, C, G
        circle circleECG = circle(E, C, G);
        draw(circleECG, dashed + blue);

        // Circle through D, A, E
        circle circleDAE = circle(D, A, E);
        draw(circleDAE, dotted + purple);

        // Circle through B, C, E
        circle circleBCE = circle(B, C, E);
        draw(circleBCE, dotted + purple);
    \end{asy}
\end{center}    
    \begin{center}
    \begin{tikzpicture}[scale=5, dot/.style={circle, fill, inner sep=0.5pt, outer sep=0pt}]

        \tkzDefPoint(0,0){A}
        \tkzDefPoint(1,0){B}
        \tkzDefPoint(1,1){C}
        \tkzDefPoint(0,1){D}
        \tkzDefPoint(0.5,0.5){O}
        \tkzDefPoint(0.7,1){E}

        % F kaki tegak lurus B ke AE; G kaki tegak lurus A ke EB
        \tkzDefPointBy[projection=onto A--E](B)\tkzGetPoint{F}
        \tkzDefPointBy[projection=onto E--B](A)\tkzGetPoint{G}
        \tkzInterLL(D,F)(C,G)\tkzGetPoint{H}
        \tkzInterLL(A,G)(B,F)\tkzGetPoint{T}

        \tkzDrawPolygon(A,B,C,D)
        \tkzDrawSegment(A,E)
        \tkzDrawSegment(B,F)
        \tkzDrawSegment(A,G)
        \tkzDrawSegment(D,F)
        \tkzDrawSegment(C,G)
        \tkzDrawSegment(E,B)
        \tkzDrawSegments(F,H H,G)
        \tkzDrawSegments(A,H H,B)
        \tkzDrawSegments[red](D,B A,C O,E O,F)
        \tkzDrawSegments[dotted,red](F,C G,D)

        \tkzDefTriangleCenter[circum](A,B,G)\tkzGetPoint{Oabg}\tkzDrawCircle(Oabg,A)
        \tkzDefTriangleCenter[circum](D,F,E)\tkzGetPoint{Odfe}\tkzDrawCircle[dashed,blue](Odfe,D)
        \tkzDefTriangleCenter[circum](E,C,G)\tkzGetPoint{Oecg}\tkzDrawCircle[dashed,blue](Oecg,E)
        \tkzDefTriangleCenter[circum](D,A,E)\tkzGetPoint{Odae}\tkzDrawCircle[dotted,purple](Odae,D)
        \tkzDefTriangleCenter[circum](B,C,E)\tkzGetPoint{Obce}\tkzDrawCircle[dotted,purple](Obce,B)

        \tkzMarkRightAngle[size=0.03](B,F,A)
        \tkzMarkRightAngle[size=0.03](A,G,E)
        \tkzMarkRightAngle[size=0.03](C,D,A)
        \tkzMarkRightAngle[size=0.03](B,C,D)
        \tkzMarkRightAngle[red, size=0.03](C,O,D)

        \tkzDrawPoints[fill=black, size=3](A,B,C,D,E,F,G,H,O,T)
        \tkzLabelPoint[below left](A){$A$}
        \tkzLabelPoint[below right](B){$B$}
        \tkzLabelPoint[above right](C){$C$}
        \tkzLabelPoint[above left](D){$D$}
        \tkzLabelPoint[above](E){$E$}
        \tkzLabelPoint[above](F){$F$}
        \tkzLabelPoint[above right](G){$G$}
        \tkzLabelPoint[below](H){$H$}
        \tkzLabelPoint[below](O){$O$}
        \tkzLabelPoint[below](T){$T$}

    \end{tikzpicture}
\end{center}

    Notasikan $\dangle$ sebagai directed angle.
    Misalkan $O$ adalah pusat persegi $ABCD$. Jelas bahwa $DB \perp AC$. Lalu misalkan $T$ sebagai perpotongan $BF$ dan $AG$ Kita punya beberapa lemma berikut
    \begin{lemmarev}\label{lemmaktomjan25-1}
        $A,B,G,O,F$ konsiklis.
        \begin{proof}
            $\dangle AFB = \dangle AOB = \dangle AGB = 90^\circ$ yang membuktikan $A,B,G,O,F$ konsiklis.
        \end{proof}
    \end{lemmarev}
    \begin{lemmarev}\label{lemmaktomjan25-2}
        $D,F,O,E$ konsiklis dan $E, O, G, C$ konsiklis.
        \begin{proof}
            Dengan bantuan Lemma~\ref{lemmaktomjan25-1} dan fakta bahwa $AB \parallel CD$ didapat
            \begin{align*}
                \dangle DOF = \dangle BOF = \dangle BAF = \dangle DEF
            \end{align*}
            yang membuktikan $D,F,O,E$ konsiklis. Analog cara tersebut, $E, O, G, C$ juga konsiklis.
        \end{proof}
    \end{lemmarev}
    \begin{lemmarev}\label{lemmaktomjan25-3}
        $A,G,E,D$ konsiklis dan $F,B,E,C$ konsiklis.
        \begin{proof}
            $\dangle ADE = 90^\circ = \dangle AGE$ yang membuktikan $A,G,E,D$ konsiklis. $\dangle BFE = 90^\circ = \dangle BCE$ yang membuktikan $F,B,E,C$ konsiklis.
        \end{proof}
    \end{lemmarev}
    \begin{lemmarev}\label{lemmaktomjan25-4}
        $F,O,G,H$ konsiklis.
        \begin{proof}
            Dari Lemma~\ref{lemmaktomjan25-1} dan Teorema Miquel pada $\triangle DHC$, didapatkan bahwa $O$ adalah Miquel point dari $\triangle DHC$ yang berakibat $F,O,G,H$ konsiklis.
        \end{proof}
    \end{lemmarev}
    Dari Lemma~\ref{lemmaktomjan25-1} dan Lemma~\ref{lemmaktomjan25-4}, didapat $F,A,H,B,G,O$ konsiklis.

    Oleh karena itu, dengan memanfaatkan semua lemma di atas kita punya
    \begin{align*}
        \angle DFC + \angle DGC &= (\angle DFE + \angle EFC) + (\angle DGE + \angle EGC)\\
        &= (\angle HFA + \angle EBC) + (\angle DAE + \angle BGH)\\
        &= (\angle HBA + \angle EBC) + (\angle DAE + \angle BAH)\\
        &= (\angle HBA + \angle BAH) + (\angle EBC + \angle DAE)\\
        &= 90^\circ + \angle BAE\\
        &= 90^\circ + (180^\circ - \angle FTG)\\
        &= 90^\circ + (180^\circ - \angle BTA)\\
        &= 90^\circ + (\angle ABF + \angle GAB)\\
        &= 90^\circ + (\angle AHF + \angle GHB)\\ 
        \angle DFC + \angle DGC &= 90^\circ + (\angle AHD + \angle BHC)\\ 
    \end{align*}
\end{proof}
